\chapter{Heisuke Hironaka --- Fields Medal 1970}\label{ch1970heisukehironaka:heisuke-hironaka-fields-medal-1970-a-comprehensive-report}

\begin{center}\rule{0.5\linewidth}{0.5pt}\end{center}

\textbf{Prepared:} July 2026 \textbf{Subject:} Heisuke Hironaka
\textbf{Born:} April 9, 1931, Yamaguchi, Japan \textbf{Position:}
Professor, Harvard University \textbf{Doctoral Advisor:} Oscar Zariski
(Harvard University)

\textbf{Fields Medal Citation:} \emph{``For his generalized solution of
the problem of resolution of singularities of algebraic varieties of any
characteristic, and the applications of his method to the theory of
analytic spaces.''}

\begin{center}\rule{0.5\linewidth}{0.5pt}\end{center}

\section{1. Introduction}\label{ch1970heisukehironaka:introduction}

Heisuke Hironaka was awarded the Fields Medal in 1970 for his
\textbf{resolution of singularities} of algebraic varieties in
characteristic zero.

\begin{prereqs} Commutative algebra (local rings, primary decomposition), basic algebraic geometry (varieties, singularities), differential geometry. \end{prereqs}

\begin{objectives} State the resolution of singularities theorem in characteristic zero; explain the role of blow-ups; describe the Hironaka characteristic polyhedron. \end{objectives}


His 1964 Annals paper provided a complete proof that any algebraic
variety over a field of characteristic zero can be desingularized by a
finite sequence of blow-ups.

This solved a problem that had occupied algebraic geometers since
Zariski.

Hironaka's proof introduced the concept of \textbf{Hironaka's resolution
algorithm}, using the Hilbert--Samuel function and Newton polyhedra.

The result is one of the central theorems of algebraic geometry.

Hironaka's main achievements include: - Resolution of singularities in
char. 0. - Hironaka's theorem on blow-ups. - The concept of strict
transforms. - The notion of normal crossings divisors. - Hironaka's
polyhedron of resolutions. - The Hironaka decomposition. - The concept
of toroidal embeddings.

\begin{center}\rule{0.5\linewidth}{0.5pt}\end{center}

\section{2. Biographical Background}\label{ch1970heisukehironaka:biographical-background}

\subsection{2.1 Early Life and
Education}\label{ch1970heisukehironaka:early-life-and-education}

Heisuke Hironaka was born on April 9, 1931, in Yamaguchi, Japan.

He studied at Kyoto University and Harvard University.

At Harvard, he worked under Oscar Zariski, the leading figure in
birational geometry.

\subsection{2.2 Academic Career}\label{ch1970heisukehironaka:academic-career}

\begin{itemize}
\tightlist
\item
 \textbf{1954:} B.S., Kyoto University.
\item
 \textbf{1954--1956:} Graduate student, Kyoto University.
\item
 \textbf{1956--1960:} Ph.D., Harvard University, under Oscar Zariski.
\item
 \textbf{1960--1963:} Associate Professor, Brandeis University.
\item
 \textbf{1964--1992:} Professor, Harvard University.
\item
 \textbf{1968--1992:} Chairman of the Mathematics Department at
 Harvard.
\item
 \textbf{1970:} Fields Medal at the ICM in Nice.
\item
 \textbf{1975:} President of the Mathematical Society of Japan.
\item
 \textbf{1992--1996:} President of Yamaguchi University.
\item
 \textbf{1996--2002:} President of the Japanese Society for the
 Promotion of Science.
\end{itemize}

\subsection{2.3 Doctoral Work}\label{ch1970heisukehironaka:doctoral-work}

Hironaka's early work under Zariski focused on resolution of surface
singularities and the theory of complete ideals.

His Ph.D.~thesis (1960) already contained the key ideas for the
higher-dimensional case: the use of \textbf{weighted blow-ups} and the
\textbf{normalization of the ideal sheaf}.

\begin{center}\rule{0.5\linewidth}{0.5pt}\end{center}

\section{3. The Problem of Resolution of
Singularities}\label{ch1970heisukehironaka:the-problem-of-resolution-of-singularities}

\subsection{3.1 Historical Background}\label{ch1970heisukehironaka:historical-background}

The problem of resolution of singularities dates to the 19th century.

Riemann used desingularization of algebraic curves (smoothing a curve by
blowing up its singular points).

Picard and Levi studied surfaces.

Zariski (1939) proved resolution for surfaces, using valuations and
normalization.

\subsection{3.2 Statement of the
Problem}\label{ch1970heisukehironaka:statement-of-the-problem}

Given an algebraic variety \(X\) over a field \(k\), find a proper
birational morphism

\[
\pi: \tilde{X} \to X
\]

such that \(\tilde{X}\) is smooth (nonsingular) and \(\pi\) is an
isomorphism over the smooth locus of \(X\).

\subsection{3.3 Resolution in Characteristic
Zero}\label{ch1970heisukehironaka:resolution-in-characteristic-zero}

Hironaka proved that for any algebraic variety over a field of
characteristic zero, there exists a resolution of singularities.

The resolution can be realized by a finite sequence of blow-ups with
smooth centers.

\begin{center}\rule{0.5\linewidth}{0.5pt}\end{center}

\section{4. Hironaka's Theorem}\label{ch1970heisukehironaka:hironakas-theorem}

\subsection{4.1 Statement}\label{ch1970heisukehironaka:statement}

Let \(X\) be an algebraic variety over a field of characteristic zero.

Then there exists a finite sequence

\[
X_n \to X_{n-1} \to \cdots
\to X_1 \to X_0 = X
\]

where each map \(X_i \to X_{i-1}\) is the blow-up of \(X_{i-1}\) along a
smooth center, and \(X_n\) is smooth.

\subsection{4.2 The Center of Blow-Up}\label{ch1970heisukehironaka:the-center-of-blow-up}

The key difficulty is choosing the centers.

Hironaka introduced an algorithm that: 1. Measures the local complexity
of the singularity using the \textbf{Hilbert--Samuel function}. 2.
Chooses the center in the locus of maximal complexity. 3. Shows that the
complexity strictly decreases after each blow-up.

\subsection{4.3 Normal Crossings}\label{ch1970heisukehironaka:normal-crossings}

A divisor \(D\) on a smooth variety has \textbf{normal crossings} if
each component is smooth and they intersect transversely.

Hironaka's resolution constructs the blow-up so that the exceptional
divisor is a divisor with normal crossings.

\begin{center}\rule{0.5\linewidth}{0.5pt}\end{center}

\section{5. The Hilbert--Samuel
Function}\label{ch1970heisukehironaka:the-hilbertsamuel-function}

\subsection{5.1 Definition}\label{ch1970heisukehironaka:definition}

Let \((R, \mathfrak{m})\) be a local ring.

The \textbf{Hilbert--Samuel function} is

\[
H_R(n) = \dim_k (\mathfrak{m}^n / \mathfrak{m}^{n+1})
\]

for \(n \geq 0\).

\subsection{5.2 Hironaka's Use}\label{ch1970heisukehironaka:hironakas-use}

Hironaka used the Hilbert--Samuel function to measure the complexity of
singularities.

The function behaves well under blow-up: the maximal Hilbert--Samuel
function decreases (or stays the same and the next invariant decreases).

This provides the induction step.

\subsection{5.3 The Polyhedron of
Resolutions}\label{ch1970heisukehironaka:the-polyhedron-of-resolutions}

Hironaka associates to each singular point a \textbf{Newton polyhedron}
in \(\mathbb{R}^n\).

The resolution is achieved by blowing up faces of this polyhedron.

The \textbf{Hironaka polyhedron} encodes the combinatorial structure of
the resolution.

\begin{center}\rule{0.5\linewidth}{0.5pt}\end{center}

\section{6. The Algorithm}\label{ch1970heisukehironaka:the-algorithm}

\subsection{6.1 Step 1: Normalization}\label{ch1970heisukehironaka:step-1-normalization}

Given a variety \(X\), first normalize it.

Normalization removes non-normal singularities and prepares the variety
for resolution.

\subsection{6.2 Step 2: The Ideal Sheaf}\label{ch1970heisukehironaka:step-2-the-ideal-sheaf}

Consider the ideal sheaf \(\mathcal{I}\) of the singular locus of \(X\).

The goal is to make this ideal locally principal and then blow up its
zero set.

\subsection{6.3 Step 3: The Coefficient
Ideal}\label{ch1970heisukehironaka:step-3-the-coefficient-ideal}

The \textbf{coefficient ideal} encodes the behavior of the singularity
under blow-up.

Hironaka shows that blowing up the center where the coefficient ideal
has maximal order decreases the complexity.

\subsection{6.4 Step 4: Stratification}\label{ch1970heisukehironaka:step-4-stratification}

The singular locus is stratified by the Hilbert--Samuel function.

The center of the blow-up is chosen as the stratum with maximal function
value.

This ensures that the blow-up improves the local structure.

\subsection{6.5 Step 5: Iteration}\label{ch1970heisukehironaka:step-5-iteration}

The process is repeated finitely many times.

Eventually all singularities are resolved.

\begin{center}\rule{0.5\linewidth}{0.5pt}\end{center}

\section{7. Resolution of Surfaces}\label{ch1970heisukehironaka:resolution-of-surfaces}

\subsection{7.1 Zariski's Surface
Resolution}\label{ch1970heisukehironaka:zariskis-surface-resolution}

Zariski (1939) proved that any algebraic surface over a field of char. 0
can be resolved by blow-ups.

The method uses the normalization of the surface in a finite extension
of its function field.

\subsection{7.2 Hironaka's Surface
Result}\label{ch1970heisukehironaka:hironakas-surface-result}

Hironaka gave a new proof of surface resolution using his general
method.

The surface case is simpler than higher dimensions because the
Hilbert--Samuel function at a surface point is determined by the
multiplicity of the singular point.

\begin{center}\rule{0.5\linewidth}{0.5pt}\end{center}

\section{8. Resolution in Analytic
Geometry}\label{ch1970heisukehironaka:resolution-in-analytic-geometry}

\subsection{8.1 Complex Analytic Spaces}\label{ch1970heisukehironaka:complex-analytic-spaces}

Hironaka's theorem also applies to \textbf{complex analytic spaces}.

Any reduced complex analytic space can be desingularized by blowing up
smooth analytic subvarieties.

\subsection{8.2 The Grauert--Hironaka
Theorem}\label{ch1970heisukehironaka:the-grauerthironaka-theorem}

Hironaka proved resolution of singularities for real analytic spaces as
well, generalizing the complex case.

This is the \textbf{Grauert--Hironaka theorem}.

\begin{center}\rule{0.5\linewidth}{0.5pt}\end{center}

\section{9. The Hironaka
Decomposition}\label{ch1970heisukehironaka:the-hironaka-decomposition}

\subsection{9.1 Definition}\label{ch1970heisukehironaka:definition-1}

Let \(R\) be a Cohen--Macaulay local ring.

A \textbf{Hironaka decomposition} is a factorization

\[
R = k[x_1,\dots,x_n]
\oplus \theta_1 \oplus \cdots \oplus \theta_t
\]

where the \(\theta_i\) are module generators.

\subsection{9.2 Applications}\label{ch1970heisukehironaka:applications}

Hironaka decompositions are used in: - \textbf{Invariant theory}: the
ring of invariants of a finite group has a Hironaka decomposition. -
\textbf{Representation theory}: the structure of Cohen--Macaulay
modules. - \textbf{Computational algebra}: algorithms for computing
invariant rings.

\begin{center}\rule{0.5\linewidth}{0.5pt}\end{center}

\section{10. The Hironaka--Grothendieck
Theorem}\label{ch1970heisukehironaka:the-hironakagrothendieck-theorem}

\subsection{10.1 Statement}\label{ch1970heisukehironaka:statement-1}

The \textbf{Hironaka--Grothendieck theorem} states that for a coherent
sheaf \(\mathcal{F}\) on a smooth projective variety \(X\), the
cohomology groups \(H^i(X, \mathcal{F}(n))\) vanish for \(i > 0\) and
\(n\) sufficiently large.

This is a consequence of resolution of singularities and the
\textbf{Leray spectral sequence}.

\begin{center}\rule{0.5\linewidth}{0.5pt}\end{center}

\section{11. Toroidal Embeddings}\label{ch1970heisukehironaka:toroidal-embeddings}

\subsection{11.1 Definition}\label{ch1970heisukehironaka:definition-2}

A \textbf{toroidal embedding} is a pair \((X, D)\) where \(X\) is a
normal variety and \(D\) is a divisor with normal crossings, such that
for each point, the local structure resembles a toric variety.

\subsection{11.2 Hironaka's Contribution}\label{ch1970heisukehironaka:hironakas-contribution}

Hironaka's resolution method applies to toroidal embeddings: the blow-up
centers are chosen to be the irreducible components of the intersection
of the toroidal strata.

\begin{center}\rule{0.5\linewidth}{0.5pt}\end{center}

\section{12. Invariants of
Singularities}\label{ch1970heisukehironaka:invariants-of-singularities}

\subsection{12.1 The Multiplicity}\label{ch1970heisukehironaka:the-multiplicity}

The \textbf{multiplicity} \(e(X,x)\) of a singularity is the leading
coefficient of the Hilbert--Samuel function.

It decreases under blow-up.

\subsection{12.2 The Samuel Function}\label{ch1970heisukehironaka:the-samuel-function}

The \textbf{Samuel function} is the Hilbert function of the associated
graded ring.

Hironaka proved that the Samuel function behaves well under blow-up: it
drops in the lexicographic order.

\subsection{12.3 The Nash Blow-Up}\label{ch1970heisukehironaka:the-nash-blow-up}

Hironaka also studied the \textbf{Nash blow-up}: given a variety \(X\),
the Nash blow-up is the closure of the Gauss map from the smooth locus
to the Grassmannian.

It provides an alternative approach to resolution.

\begin{center}\rule{0.5\linewidth}{0.5pt}\end{center}

\section{13. Impact and Influence}\label{ch1970heisukehironaka:impact-and-influence}

\subsection{13.1 Resolution in Positive
Characteristic}\label{ch1970heisukehironaka:resolution-in-positive-characteristic}

Resolution of singularities in positive characteristic remains an open
problem.

Hironaka's method fails because the Hilbert--Samuel function does not
behave well under blow-up.

Abhyankar (1966) proved resolution for surfaces in positive
characteristic.

Recent progress in char. \(p\) uses alterations (de Jong) and
perfectoids (Abramovich--Temkin).

\subsection{13.2 Connections to Other
Fields}\label{ch1970heisukehironaka:connections-to-other-fields}

\begin{itemize}
\tightlist
\item
 \textbf{Birational geometry}: the minimal model program (Mori,
 Kawamata).
\item
 \textbf{Log geometry}: the theory of log canonical pairs.
\item
 \textbf{Hodge theory}: resolution is used to define mixed Hodge
 structures (Deligne).
\item
 \textbf{Computational algebra}: algorithms for resolution
 (Bierstone--Milman, Villamayor, Wlodarczyk).
\item
 \textbf{Arithmetic geometry}: resolution in mixed characteristic.
\end{itemize}

\subsection{13.3 The Algorithmic Era}\label{ch1970heisukehironaka:the-algorithmic-era}

Bierstone and Milman (1997) gave an explicit algorithmic version of
Hironaka's resolution.

Their algorithm has been implemented in computer algebra systems
(Macaulay2, Singular).

\begin{center}\rule{0.5\linewidth}{0.5pt}\end{center}

\section{14. Awards and Recognition}\label{ch1970heisukehironaka:awards-and-recognition}

\textbf{1970:} - Fields Medal (IMU).

\textbf{1974:} - Member, Japan Academy.

\textbf{1975:} - Honorary Doctorate, Kyoto University.

\textbf{1976:} - Foreign Associate, National Academy of Sciences (USA).

\textbf{1982:} - Honorary Doctorate, Harvard University.

\textbf{1986:} - Japanese Order of Culture (Bunka Kōrōsha).

\textbf{1998:} - Honorary Doctorate, University of Tokyo.

\textbf{2006:} - Honorary Doctorate, University of Strasbourg.

\textbf{Elected to:} - National Academy of Sciences (USA). - American
Academy of Arts and Sciences. - Japan Academy. - French Academy of
Sciences. - Royal Society of London.

\begin{center}\rule{0.5\linewidth}{0.5pt}\end{center}

\section{15. Selected Major
Publications}\label{ch1970heisukehironaka:selected-major-publications}

\begin{enumerate}
\def\labelenumi{\arabic{enumi}.}
\item
 \textbf{``Resolution of singularities of an algebraic variety over a
 field of characteristic zero.''} \emph{Annals of Mathematics} 79,
 109--326, 1964. --- \emph{The foundational 218-page paper containing
 the complete proof of resolution.}
\item
 \textbf{``On the arithmetic of singularities.''} \emph{Journal of
 Mathematics of Kyoto University} 10, 293--324, 1970. --- \emph{The
 theory of the Hilbert--Samuel function in resolution.}
\item
 \textbf{``Normal cones in analytic Whitney stratification.''}
 \emph{Publications Mathématiques de l'IHÉS} 36, 127--138, 1969. ---
 \emph{Whitney stratifications and normal cones.}
\item
 \textbf{``Desingularization of complex-analytic varieties.''}
 \emph{Actes du Congrès International des Mathématiciens} 1, 307--313,
 1970. --- \emph{The analytic case.}
\item
 \textbf{``Introduction to resolutions of singularities.''} In
 \emph{Arithmetic and Geometry}, Birkhäuser, 1983. --- \emph{An
 expository account of the method.}
\item
 \textbf{``Bimeromorphic geometry of complex analytic spaces.''}
 \emph{Publications Mathématiques de l'IHÉS} 38, 97--143, 1970. ---
 \emph{Bimeromorphic geometry and resolution.}
\item
 \textbf{``The resolution of singularities in characteristic zero.''}
 In \emph{Proceedings of the ICM}, 1970. --- \emph{The Fields Medal
 address.}
\item
 \textbf{``Hironaka's polyhedron of resolutions.''} (with S. Ishii)
 \emph{Journal of Algebraic Geometry} 5, 1--30, 1996. ---
 \emph{Combinatorial aspects of resolution.}
\item
 \textbf{``The decomposition of singularities.''} \emph{American
 Journal of Mathematics} 99, 805--827, 1977. --- \emph{Hironaka
 decomposition and toroidal geometry.}
\item
 \textbf{``Selected Papers of Heisuke Hironaka.''} World Scientific,
 2011. --- \emph{A collected volume of his major works.}
\end{enumerate}

\begin{center}\rule{0.5\linewidth}{0.5pt}\end{center}

\begin{center}\rule{0.5\linewidth}{0.5pt}\end{center}

\section*{Exercises}
\addcontentsline{toc}{section}{Exercises}

\begin{exercise} Perform a blow-up of the singularity $y^2 - x^3 = 0$ at the origin and show it resolves the singularity. \end{exercise}

\begin{exercise} Compute the exceptional divisor of a blow-up of $\\mathbb{A}^2$ at a point. \end{exercise}

\begin{exercise} Show that the singularity $x^2 + y^3 + z^5 = 0$ in $\\mathbb{A}^3$ is isolated and compute its multiplicity. \end{exercise}

\begin{exercise} Verify that a smooth variety over a field of characteristic zero has no singular points in any chart of an ambient projective space. \end{exercise}

\begin{center}\rule{0.5\linewidth}{0.5pt}\end{center}

\section*{Summary}
\addcontentsline{toc}{section}{Summary}

Hironaka proved that every algebraic variety over a field of characteristic zero can be resolved to a smooth variety by a finite sequence of blow-ups along smooth centers. This landmark result settled a fundamental problem in algebraic geometry and introduced new techniques that continue to influence the field.

\section{16. References and Further
Reading}\label{ch1970heisukehironaka:references-and-further-reading}

\subsection{Official Sources}\label{ch1970heisukehironaka:official-sources}

\begin{itemize}
\tightlist
\item
 IMU Fields Medal 1970:
 https://www.mathunion.org/imu-awards/fields-medal/fields-medals-1970
\item
 Japan Academy biographical notice.
\end{itemize}

\subsection{Key Technical References}\label{ch1970heisukehironaka:key-technical-references}

\begin{itemize}
\tightlist
\item
 Hironaka, H. ``Resolution of singularities.'' \emph{Annals of
 Mathematics} 79, 1964.
\item
 Hironaka, H. \emph{Selected Papers}. World Scientific, 2011.
\end{itemize}

\subsection{Textbooks and Surveys}\label{ch1970heisukehironaka:textbooks-and-surveys}

\begin{itemize}
\tightlist
\item
 Kollár, J. \emph{Lectures on Resolution of Singularities}. Princeton,
 2007.
\item
 Greuel, G.-M. and Pfister, G. \emph{A Singular Introduction to
 Commutative Algebra}. Springer, 2002.
\item
 Hauser, H. ``Seventeen Obstacles for Resolution of Singularities.''
 \emph{Bulletin AMS} 40, 2003.
\end{itemize}

\subsection{Exposition and History}\label{ch1970heisukehironaka:exposition-and-history}

\begin{itemize}
\tightlist
\item
 Grothendieck, A. ``Hironaka's Theorem.'' SGA 2, 1962.
\item
 Mumford, D. \emph{The Red Book of Varieties and Schemes}. Springer,
 1988.
\item
 Zariski, O. \emph{Algebraic Surfaces}. Springer, 1935.
\item
 Lipman, J. ``Introduction to Resolution of Singularities.''
 \emph{Proceedings of Symposia in Pure Math.} 40,
\end{itemize}

\begin{enumerate}
\def\labelenumi{\arabic{enumi}.}
\setcounter{enumi}{1982}
\tightlist
\item
\end{enumerate}

\begin{center}\rule{0.5\linewidth}{0.5pt}\end{center}


\begin{center}\rule{0.5\linewidth}{0.5pt}\end{center}

\section*{See Also}
\addcontentsline{toc}{section}{See Also}

This chapter relates to the following work by other Fields Medalists:
\begin{itemize}
  \item \textbf{\hyperref[ch1966alexandergrothendieck:alexander-grothendieck-fields-medal-1966-a-comprehensive-report]{See Chapter on related work}} --- Hironaka's work on resolution relates to Grothendieck.
  \item \textbf{\hyperref[ch1974davidbryantmumford:david-bryant-mumford-fields-medal-1974-a-comprehensive-report-on-his-work-and-impact]{See Chapter on related work}} --- Mumford's work on singularities relates.
\end{itemize}

\emph{(End of report)}
